%=====================================================================
\chapter{Major Hypervisor Platforms}\label{ch:platforms}
%=====================================================================
\locator{2}{Part II --- Hypervisors and Compute Virtualization}

\begin{outcomes}
By the end of this chapter you will be able to:
\begin{tight}
  \item \textbf{Describe} the architecture of ESXi, Hyper-V, KVM and Xen, and
        \textbf{say} for each one where the device drivers live.
  \item \textbf{Explain} why the location of the drivers is the single fact
        that most determines a platform's security and performance character.
  \item \textbf{Compare} the four on licensing, ecosystem, guest support and
        operational tooling.
  \item \textbf{Compute} the three-year cost of a small cluster on two of them,
        and \textbf{identify} which cost dominates.
  \item \textbf{Recommend} a platform for a stated requirement and defend it
        against the strongest alternative.
\end{tight}
\end{outcomes}

\begin{prereq}
\chref{architectures} (Type-1 and Type-2, and the three techniques) and
\chref{hardware} (VT-x, EPT, the IOMMU). Every platform here is a different set
of engineering answers built on the same silicon, so the hardware chapter is
what lets you see what they actually differ about.
\end{prereq}

\begin{keyterms}
ESXi $\bullet$ VMkernel $\bullet$ vSphere $\bullet$ vCenter $\bullet$ Hyper-V
$\bullet$ parent partition $\bullet$ child partition $\bullet$ VMBus $\bullet$
enlightenment $\bullet$ KVM $\bullet$ QEMU $\bullet$ libvirt $\bullet$ Xen
$\bullet$ dom0 $\bullet$ domU $\bullet$ PV $\bullet$ HVM $\bullet$ PVH
$\bullet$ disaggregation $\bullet$ driver domain $\bullet$ trusted computing
base
\end{keyterms}

\section{The Question That Separates Them}

All four platforms in this chapter run on the same processors, use VT-x or
AMD-V the same way, and use extended page tables for memory. At the level of
\chref{hardware} they are nearly identical.

What separates them is a question \chref{hardware} did not ask: \emph{where do
the device drivers live?} A hypervisor needs drivers for every disk controller
and network card in the machine, and drivers are large, written by hardware
vendors, and the commonest source of kernel bugs. Any code that can crash the
host is part of the \term{trusted computing base}, and a driver can always
crash the host.

The four platforms give four different answers, and almost every other
difference between them follows.

\dsfig{%
\begin{tikzpicture}[font=\scriptsize,
  box/.style={rounded corners=2pt, line width=0.8pt, align=center,
              inner sep=2pt, minimum height=0.5cm, font=\scriptsize},
  hw/.style={box, draw=neutral, fill=neutral!15, text=neutral!40!black,
             minimum width=3.3cm},
  hv/.style={box, draw=primary, fill=primary!14, text=primary,
             font=\scriptsize\bfseries, minimum width=3.3cm},
  gu/.style={box, draw=secondary, fill=secondary!8, text=secondary!80!black,
             minimum width=1.5cm},
  pr/.style={box, draw=purplehl, fill=purplehl!12, text=purplehl!70!black,
             minimum width=1.5cm},
  dr/.style={box, draw=accent, fill=accent!12, text=accent!70!black,
             font=\tiny, minimum width=1.5cm, minimum height=0.4cm},
  ttl/.style={font=\scriptsize\bfseries, text=primary},
  blurb/.style={font=\tiny\itshape, text=neutral, align=center}
]
% ---------------- ESXi ----------------
\node[ttl] at (1.9,3.5) {ESXi};
\node[hw] (e0) at (1.9,0)   {hardware};
\node[hv] (e1) at (1.9,0.75) {VMkernel};
\node[dr, minimum width=3.1cm] (ed) at (1.9,1.3) {drivers inside it};
\node[gu] (e2) at (1.1,2.2) {guest}; \node[gu] (e3) at (2.7,2.2) {guest};
\node[blurb] at (1.9,-0.75) {one purpose-built kernel\\of about 150 MB};

% ---------------- Hyper-V ----------------
\node[ttl] at (7.1,3.5) {Hyper-V};
\node[hw] (h0) at (7.1,0)   {hardware};
\node[hv] (h1) at (7.1,0.75) {hypervisor (thin)};
\node[pr, minimum width=1.6cm] (hp) at (6.0,1.75) {parent\\partition};
\node[dr, minimum width=1.6cm] (hd) at (6.0,2.45) {Windows\\drivers};
\node[gu] (h2) at (7.9,1.75) {child}; \node[gu] (h3) at (7.9,2.45) {child};
\draw[-{Stealth[length=1.4mm]}, accent!70, line width=0.7pt] (7.15,1.9) -- (6.8,1.9);
\node[font=\tiny, text=accent!70!black] at (7.35,2.1) {VMBus};
\node[blurb] at (7.1,-0.75) {thin hypervisor; drivers\\in a privileged Windows};

% ---------------- KVM ----------------
\node[ttl] at (12.3,3.5) {KVM};
\node[hw] (k0) at (12.3,0)   {hardware};
\node[hv] (k1) at (12.3,0.9) {Linux kernel\\{\tiny\mdseries + kvm.ko}};
\node[dr, minimum width=3.1cm] (kd) at (12.3,1.6) {every Linux driver};
\node[gu, minimum width=1.5cm] (k2) at (11.5,2.5) {QEMU\\+ guest};
\node[gu, minimum width=1.5cm] (k3) at (13.1,2.5) {QEMU\\+ guest};
\node[blurb] at (12.3,-0.75) {a general-purpose kernel\\that is also the
                            hypervisor};

% ---------------- Xen ----------------
\node[ttl] at (17.5,3.5) {Xen};
\node[hw] (x0) at (17.5,0)   {hardware};
\node[hv] (x1) at (17.5,0.75) {Xen hypervisor};
\node[pr, minimum width=1.6cm] (xp) at (16.4,1.75) {dom0};
\node[dr, minimum width=1.6cm] (xd) at (16.4,2.45) {Linux\\drivers};
\node[gu] (x2) at (18.3,1.75) {domU}; \node[gu] (x3) at (18.3,2.45) {domU};
\draw[-{Stealth[length=1.4mm]}, accent!70, line width=0.7pt] (17.55,1.9) -- (17.2,1.9);
\node[font=\tiny, text=accent!70!black] at (17.8,2.1) {rings};
\node[blurb] at (17.5,-0.75) {hypervisor of about 1 MB;\\drivers in a
                            privileged guest};

\foreach \a/\b in {e1/e0, h1/h0, k1/k0, x1/x0}{
  \draw[-{Stealth[length=1.4mm]}, primary!55, line width=0.7pt] (\a) -- (\b);}
\end{tikzpicture}}%
{The four platforms, drawn to the same question: where are the drivers (red)?
Xen's answer gives it the smallest hypervisor in the room and the most moving
parts; ESXi's gives it the most predictable behaviour and the shortest hardware
compatibility list.}{four-platforms}

\section{VMware ESXi and vSphere}

ESXi is a purpose-built Type-1 hypervisor. Its kernel, the \term{VMkernel}, was
written for one job and contains the scheduler, the memory manager and the
device drivers. Nothing general-purpose runs beside it; the shell is
deliberately limited and disabled by default.

\term{vSphere} is the product that matters commercially: ESXi hosts plus
\term{vCenter}, the management server that turns a set of hosts into a cluster
with shared storage, automated placement, live migration and failover. Almost
everything an enterprise buys VMware for lives in vCenter rather than in ESXi.

\begin{conceptbox}[The Compatibility List Is the Point, Not a Limitation]
ESXi runs on a short list of certified hardware, and students usually read this
as a weakness. It is the central design decision.

\smallskip
Because VMware controls which drivers exist, it can test them, and because the
VMkernel does nothing else, its behaviour under load is predictable in a way a
general-purpose kernel's is not. An estate of certified hardware running a
tested driver set is why vSphere acquired a reputation for not surprising
people at three in the morning --- and that reputation, rather than any
technical feature, is what the licence buys.

\smallskip
The cost is real: a host that is not on the list may not install at all, and
a new network card may wait months for a driver.
\end{conceptbox}

\section{Microsoft Hyper-V}

Hyper-V is also Type-1, and the picture most people hold of it is wrong.
Installing the Hyper-V role on Windows Server does not make Windows the
hypervisor; it slides a thin hypervisor \emph{underneath} the running Windows,
which is then demoted to the \term{parent partition} --- a privileged guest.

The parent partition owns the physical devices and runs the Windows driver
stack. \term{Child partitions} --- the ordinary guests --- have no direct device
access and talk to the parent over \term{VMBus}, a shared-memory channel. A
guest with Hyper-V's integration services installed uses \term{enlightened}
paravirtualized drivers over VMBus; a guest without them falls back to emulated
devices and the performance of \chref{architectures}'s worked example.

\begin{alertbox}[Hyper-V changes the machine underneath you]
Once the Hyper-V role is enabled, the Windows you were using is a guest. This
has consequences people meet by accident:

\smallskip
Other hypervisors stop working, or quietly become nested. VirtualBox and VMware
Workstation on a Hyper-V-enabled Windows run on top of Hyper-V's platform
rather than on the hardware, and historically simply refused. Benchmarks taken
on that machine are measuring a nested guest.

\smallskip
It is also enabled by features that do not mention it: Windows Subsystem for
Linux 2, Windows Sandbox, Docker Desktop and memory-integrity protection all
turn Hyper-V on. A great deal of ``my virtual machines got slow after a Windows
update'' is this.
\end{alertbox}

\section{KVM and QEMU}

KVM is a kernel module. Loading \texttt{kvm.ko} and \texttt{kvm\_intel.ko} turns
Linux into a hypervisor by giving it the ability to enter VMX non-root
operation; the module is about twenty thousand lines, because the scheduler,
memory manager and drivers it needs are already there.

That is also the whole architectural claim: \textbf{a guest is a process}. It
is scheduled by the ordinary Linux scheduler, its memory is ordinary process
memory, \texttt{top} shows it, \texttt{kill} ends it, cgroups limit it. Every
tool that works on a Linux process works on a virtual machine.

KVM provides only the processor and memory virtualization. The virtual
hardware --- disk controllers, network cards, graphics, BIOS --- is QEMU's job,
running in user space as that process. \term{libvirt} sits above both and
provides the stable API that \texttt{virsh}, \texttt{virt-manager}, OpenStack
and oVirt all drive.

\begin{conceptbox}[One Design Decision, Two Opposite Consequences]
\textbf{For:} KVM inherits everything. Every driver Linux has, every filesystem,
every scheduler improvement, every networking feature, with no porting work.
When a new network card ships with a Linux driver, KVM supports it that day.
This is why KVM runs most of the world's public cloud.

\smallskip
\textbf{Against:} KVM inherits everything. The trusted computing base is the
entire Linux kernel --- millions of lines, every driver loaded, every
filesystem --- and a bug in any of it is a host compromise. ESXi's 150 MB and
Xen's 1 MB are the competing answers to exactly this.

\smallskip
Neither consequence can be had without the other, and which one dominates
depends on whether you are optimising for hardware support or for attack
surface.
\end{conceptbox}

\section{Xen}

Xen is the smallest hypervisor here --- on the order of a megabyte --- because
it deliberately contains almost nothing. No drivers, no filesystems, no network
stack. It schedules domains, manages memory, and passes messages.

Everything else lives in \term{dom0}, a privileged guest (normally Linux) that
owns the physical devices and runs the driver stack. Unprivileged guests are
\term{domU} and reach devices through split drivers terminating in dom0.

Xen was the paravirtualization project of \chref{vmm}, and its guest modes
record that history: \term{PV} (fully paravirtualized, no hardware assistance
needed, guest kernel modified), \term{HVM} (hardware-assisted, unmodified
guests, from 2006), and \term{PVH} (hardware assistance for the processor,
paravirtualized devices, no emulated chipset --- the modern default).

\begin{conceptbox}[Disaggregation: the idea Xen got right early]
Because dom0 is a guest rather than part of the hypervisor, it can be split
up. A \term{driver domain} is an unprivileged domain given one device via the
IOMMU, running that device's driver and nothing else. A buggy or hostile
network driver then crashes a domain that can be restarted, instead of the
host.

\smallskip
This is \term{disaggregation}, and it is the strongest security argument any
platform in this chapter can make: it is the only one that can reduce the
blast radius of a driver bug to something smaller than the whole machine. It
is also operationally fiddly, which is why it is common in security-focused
deployments such as Qubes OS and rare elsewhere.

\smallskip
Note where the idea goes next. \chref{lightweight} removes device emulation
from the trusted path almost entirely, and \chref{gpu} uses the same IOMMU
machinery to hand accelerators to guests. Xen's answer was early rather than
wrong.
\end{conceptbox}

\section{Choosing}

\rowcolors{2}{white}{lightbg}
\begin{longtable}{@{}L{2.5cm}L{3.0cm}L{3.0cm}L{3.0cm}L{3.0cm}@{}}
\toprule
\rowcolor{primary!15}
& \textbf{ESXi} & \textbf{Hyper-V} & \textbf{KVM} & \textbf{Xen} \\
\midrule
\endhead
Hypervisor size & $\approx$150 MB & thin, plus a full Windows parent &
the Linux kernel & $\approx$1 MB, plus dom0 \\
Drivers live in & the VMkernel & the parent partition & the Linux kernel &
dom0 or a driver domain \\
Licence & Commercial; free tier withdrawn in 2024 & Included with Windows
Server; free Hyper-V Server discontinued & GPL, free & GPL, free \\
Guest support & Very broad, certified & Windows strongest; Linux good &
Anything x86 or ARM & Anything, three guest modes \\
Management & vCenter --- the strongest in the industry & System Center, Windows
Admin Center & libvirt, oVirt, Proxmox, OpenStack & XenCenter, XCP-ng \\
Used by & Enterprise data centres & Windows-centred estates & AWS Nitro,
Google, DigitalOcean, most OpenStack & AWS before Nitro, Citrix, Qubes \\
Strongest when & Mixed hardware, mixed guests, operations staff who must not be
surprised & The estate is already Windows and licensed & You want control, no
licence cost, and newest hardware & Minimal trusted base, or driver
isolation \\
\bottomrule
\end{longtable}

\begin{worked}[Three Years of a Small Cluster]
Size and cost the three-host cluster of \chref{what}'s worked example on vSphere
and on KVM. Each host has 2 sockets $\times$ 32 cores.

\smallskip
\textbf{vSphere.} VMware licenses per core with a 16-core minimum per socket.
Here each socket has 32 cores, so the minimum does not bite:
\[ 3 \text{ hosts} \times 2 \text{ sockets} \times 32 = 192 \text{ cores} \]
At a list price of about \$350 per core per year for a standard edition, plus
vCenter at about \$12{,}000 for three years:
\[ 192 \times 350 \times 3 + 12{,}000 = 201{,}600 + 12{,}000
   = \mathbf{\$213{,}600} \]

\smallskip
\textbf{KVM.} The software is free. A supported distribution with a virtualisation
subscription runs about \$1{,}500 per host per year:
\[ 3 \times 1{,}500 \times 3 = \mathbf{\$13{,}500} \]
Unsupported, it is \$0.

\smallskip
\textbf{Step 1 --- the gap.} \$200{,}100 over three years, or about
\$66{,}700 a year.

\smallskip
\textbf{Step 2 --- what the gap must buy.} At a loaded cost of roughly
\$45{,}000 a year for an experienced Linux engineer, the licence is worth
paying if vSphere saves about \textbf{1.5 full-time engineers} of effort
against a self-managed KVM estate --- or if it prevents outages worth that
much.

\smallskip
\textbf{Step 3 --- read it honestly in both directions.} For an estate of three
hosts the answer is almost always KVM: three hosts do not need 1.5 engineers
of management either way, so the licence buys little. For an estate of three
hundred hosts with a mixed hardware fleet, staff who know vCenter and an
availability target with contractual penalties, the same arithmetic inverts.

\smallskip
\textbf{What the arithmetic leaves out.} Migration cost. An estate already on
vSphere carries years of runbooks, automation, backup integration and staff
knowledge, and moving is a project, not a decision. Prices also change: VMware's
2024 move to subscription-only bundles raised many renewals severalfold and
sent a visible part of the market to Proxmox and OpenStack. Check current
pricing before quoting any of these figures.
\end{worked}

\begin{pitfall}
\begin{tight}
  \item \textbf{Thinking Hyper-V runs on top of Windows.} It runs underneath
        it. Once enabled, your Windows is a privileged guest --- which is why
        other hypervisors on the same machine become nested or refuse.
  \item \textbf{Describing KVM as ``just a kernel module'' and stopping
        there.} The module is the small part. QEMU supplies every virtual
        device and libvirt supplies the management API; a KVM deployment is all
        three.
  \item \textbf{Comparing hypervisor sizes without counting dom0 or the parent
        partition.} Xen's hypervisor is a megabyte and its dom0 is a full Linux
        kernel. The honest comparison is of the whole privileged path.
  \item \textbf{Assuming the free tier still exists.} VMware withdrew free ESXi
        in 2024 and Microsoft discontinued the free Hyper-V Server. Student and
        evaluation routes change; check before building a course or a budget on
        one.
  \item \textbf{Choosing a platform on hypervisor features.} All four
        virtualize a processor competently. The decision is made on management
        tooling, licence cost, hardware support and the skills of the people
        who will be on call.
\end{tight}
\end{pitfall}

\begin{lab}[The Same Guest on Two Platforms]
Builds one guest under KVM and the same guest under VirtualBox, then compares
what each platform exposes to it. Appendix~A has both. Expect forty minutes.

\begin{lstlisting}[language=bash]
#!/usr/bin/env bash
# ch05_platforms.sh -- one workload, two platforms, and what differs.
set -euo pipefail

# --- 1. What is this host already running, if anything? ----------------
systemd-detect-virt || echo "bare metal"
lsmod | grep -E '^(kvm|kvm_intel|kvm_amd|vboxdrv)' || true

# --- 2. Define a guest on KVM, from the Chapter 1 image. ---------------
sudo virsh dumpxml vss1 > vss1.xml
grep -E '<(domain|emulator|controller|interface|model)' vss1.xml | head -20
# Note the machine type: QEMU presents a specific, named virtual chipset.

# --- 3. Ask the guest what hardware it thinks it has. ------------------
sudo virsh start vss1 ; sleep 25
IP=$(sudo virsh domifaddr vss1 | awk '/ipv4/{split($4,a,"/"); print a[1]}')
G () { ssh -o StrictHostKeyChecking=no "ubuntu@$IP" "$1"; }
G "sudo dmidecode -s system-manufacturer"
G "sudo dmidecode -s system-product-name"
G "lspci | head; systemd-detect-virt"

# --- 4. Build the same guest under VirtualBox. ------------------------
VBoxManage createvm --name vss2 --ostype Ubuntu_64 --register
VBoxManage modifyvm vss2 --memory 2048 --cpus 2 --nic1 nat
VBoxManage clonehd ubuntu.img vss2.vdi --format VDI
VBoxManage storagectl vss2 --name SATA --add sata
VBoxManage storageattach vss2 --storagectl SATA --port 0 \
           --type hdd --medium vss2.vdi
VBoxManage startvm vss2 --type headless

# --- 5. Compare what the two platforms claim to be. -------------------
# Each hypervisor identifies itself: the manufacturer string, the PCI
# device list and the CPUID hypervisor leaf all differ.
echo "KVM guest reports:"      ; G "sudo dmidecode -s system-manufacturer"
echo "VirtualBox guest reports: innotek GmbH (check in the guest console)"

# --- 6. The feature that is not in the free product. ------------------
# Live migration needs two hosts and shared storage. virsh offers it;
# VirtualBox does not. Chapter 17 is where this becomes the whole point.
sudo virsh migrate --help | head -5

# --- 7. Tear down. ----------------------------------------------------
sudo virsh destroy vss1 || true
VBoxManage controlvm vss2 poweroff || true
VBoxManage unregistervm vss2 --delete || true
\end{lstlisting}

\textbf{What to notice.} Step 5 is the honest answer to ``can a guest tell?'' ---
every platform identifies itself in several places, and none of them tries very
hard to hide. Step 6 is the difference that actually decides procurement:
the hypervisors are comparable, and the operational features are not.
\end{lab}

\begin{checkpoint}
\begin{tightnum}
  \item For each of ESXi, Hyper-V, KVM and Xen, say in one phrase where the
        device drivers run, and give the consequence for the trusted computing
        base.
  \item A colleague reports that after installing Docker Desktop on Windows,
        their VirtualBox guests became much slower. Explain what happened.
  \item Using the worked example's method, at what cluster size does a vSphere
        licence begin to look defensible, and what assumption in your answer is
        doing the most work?
\end{tightnum}
\end{checkpoint}

\begin{chaptersummary}
\begin{tight}
  \item All four platforms use the same hardware features. What separates them
        is where the device drivers live, and almost every other difference
        follows from that.
  \item ESXi: a purpose-built kernel of about 150 MB with drivers inside it, a
        certified hardware list, and vCenter --- which is what the licence
        actually buys.
  \item Hyper-V: a thin hypervisor that slides underneath a running Windows,
        demoting it to a privileged parent partition that owns the drivers.
        Enabling it changes the machine, often by accident.
  \item KVM: a kernel module that makes Linux a hypervisor, so a guest is a
        process. It inherits every Linux driver and the entire Linux kernel as
        trusted base --- the same fact read twice.
  \item Xen: about a megabyte of hypervisor with nothing in it; dom0 owns the
        drivers, and driver domains can shrink the blast radius of a driver bug
        to less than the host.
  \item Licence cost differences are large --- roughly \$200{,}000 over three
        years for a three-host cluster --- and are decided by staff effort and
        risk, not by hypervisor features.
\end{tight}
\end{chaptersummary}

\begin{reviewq}
  \item \textbf{(MCQ)} In Hyper-V, the Windows installation on which the role
        was enabled becomes: (a)~the hypervisor; (b)~the parent partition, a
        privileged guest; (c)~an ordinary child partition; (d)~unchanged.
  \item \textbf{(MCQ)} Xen's hypervisor is roughly a megabyte because:
        (a)~it is written in assembly; (b)~it contains no device drivers,
        filesystems or network stack; (c)~it supports fewer guests;
        (d)~it is compiled with optimisation.
  \item \textbf{(Short)} State the single architectural claim KVM makes about
        what a guest is, and give three operational consequences of it.
  \item \textbf{(Short)} Explain why ESXi's short hardware compatibility list
        is a design decision rather than a shortcoming, and what it costs.
  \item \textbf{(Short)} Describe a driver domain and the specific failure it
        is meant to contain.
  \item \textbf{(Applied)} A department has six hosts, two Linux engineers, no
        existing virtualization and a small budget. A vendor proposes vSphere.
        Write the recommendation you would give, including the two questions
        you would ask before committing and what answer to each would change
        your mind.
  \item \textbf{(Applied)} A hosting provider must run untrusted guests and is
        choosing between KVM and Xen with driver domains. Set out the security
        argument for each, name what the Xen design costs operationally, and
        say what you would need to measure before deciding.
\end{reviewq}
